\section*{Appendix: AI prompts}

Every input to the AI system for “Walking breaks \& 100\% focus\_time”, in order, with the author's classification. Live reviews and checks the plugin started on its own are marked as such.

\begin{enumerate}
  \item \textbf{2026-10-01 10:00 · Feedback and critique of own text · adapted (Einleitung)} \\
  Gib mir Feedback zur Einleitung \& zu 50\% der Methodik
  \\ \textit{Answer (claude-opus-5-5):} Hinweis 1
  \item \textbf{2026-10-02 10:00 · Grammar, spelling and punctuation corrections · not annotated · [AI-TEXT]} \\
  Korrigiere Rechtschreibung
  \\ \textit{Answer (claude-opus-5-5):} [AI-TEXT] bla
  \item \textbf{2026-10-03 10:00 · Live review while writing (short hints on changed paragraphs) · hints only, no classification required} \\
  Live review (grammar-guy) of main.tex l.1–3
  \\ \textit{Answer (haiku):} (no hints)
  \item \textbf{2026-10-04 10:00 · Checks of field-specific terminology · not annotated} \\
  Use the cowrite:terminologist agent to review main.tex
  \\ \textit{Answer (claude-opus-5-5):} The Terminologist is on it.
  \item \textbf{2026-10-05 10:00 · Checks of field-specific terminology · in reply to 2026-10-04 10:00 · not annotated} \\
  Reply to The Terminologist (the cowrite:terminologist agent) on its terminology feedback from 10:00, log entry a4. My reply: the field says subjects.
  \\ \textit{Answer (m):} Asking.
  \item \textbf{2026-10-05 10:01 · Checks of field-specific terminology · report from terminologist · in reply to 2026-10-05 10:00 · not annotated} \\
  <agent-message from="x"> [Subagent hand-back] Fair point for psychology; in medicine it is participants. </agent-message>
  \\ \textit{Answer (m):} The Terminologist concedes for psychology.
  \item \textbf{2026-10-06 10:00 · Briefing lines for a reviewer persona, distilled from a feedback discussion; written to party.md only once the author confirmed them · adopted} \\
  Briefing proposal: what The Terminologist should remember from its input a4, for party.md
  \\ \textit{Answer (m):} - Keep "subjects": the field's term.
  \\ \textit{Comment:} - Keep "subjects": the field's term.
  \item \textbf{2026-10-07 09:00 · Structure and argumentation feedback · rejected (Methods)} \\
  A question over several lines, with a\_underscore, a \{brace\}, 100\% and a very long tail word word word word word word word word word word word word word word word word word word word word word word word word word word word word word word word word word word word word word word word word
  \\ \textit{Answer (claude-opus-5-5):} \#\# Findings | Where | What | |---|---| | main.md:12 | **unclear** `claim` | Why this order?
  \\ \textit{Comment:} Not for this\_journal: 100\% \& more
  \item \textbf{2026-10-07 09:30 · Structure and argumentation feedback · adopted (Discussion)} \\
  Check the outline
  \\ \textit{Answer (claude-opus-5-5):} Outline feedback
  \item \textbf{2026-10-07 09:40 · Check of the planned points against the text · hints only, no classification required} \\
  Scope check of the planned points
  \\ \textit{Answer (haiku):} []
  \item \textbf{2026-10-07 09:50 · Feedback and critique of own text · not annotated} \\
  Feedback on the abstract, please
  \\ \textit{Answer (claude-opus-5-5):} Three issues
\end{enumerate}
